\chapter{Lòng Tự Trọng, Xấu Hổ Và Khả Năng Đứng Dậy}

\begin{leadbox}
Rút gọn từ xấu hổ là điều khó khăn nhất mà một học sinh phải học.
\textit{(Rising from shame is one of the hardest things a student must learn.)}
Lòng tự trọng bị thương có thể thôi thúc hoặc có thể tê liệt.
\textit{(Wounded dignity can motivate or paralyze.)}
Mười tình huống sau nhìn vào những khoảnh khắc khi học sinh chọn đứng dậy.
\textit{(The 10 situations below look at moments when students choose to rise.)}
\end{leadbox}

% Story 51
\begin{truyen}{51. Đọc Bài Bị Cười}{Read Aloud and Mocked}
Em Mai được yêu cầu đọc bài đầu tiên trong buổi học.
\textit{(Mai was asked to read first in class.)}
Từ dòng đầu tiên, em đã đọc sai từng tiếng.
\textit{(From the first line, she mispronounced every word.)}

Lớp bắt đầu cười.
\textit{(The class started laughing.)}
Không phải cười với bài Mai, mà cười với những từ sai lệch của em.
\textit{(Not laughing at Mai, but at her mispronunciations.)}

Em Mai cô đỏ và dừng đọc.
\textit{(Mai turned red and stopped reading.)}
Cô giáo đưa sách cho em khác.
\textit{(The teacher passed the book to another student.)}

Từ hôm sau, em Mai không bao giờ dùng tay nữa.
\textit{(After that, Mai never raised her hand again.)}
Cô bị im lặng, không tham gia lớp, sợ bất cứ cơ hội nào để phải nói trước mọi người.
\textit{(She went silent, didn't participate, feared any chance to speak in front of the class.)}

Ba tháng sau, cô giáo mới nhận ra.
\textit{(Three months later, the teacher finally noticed.)}
Cô gọi em nói chuyện riêng.
\textit{(She called Mai to talk privately.)}
``Cô thích nghe được giọng em. Không ai hoàn toàn không làm sai lần đầu.''
\textit{(``I'd like to hear your voice. No one gets it right the first time.'')}

Em Mai bắt đầu dỡi lại từ từ.
\textit{(Mai slowly began to recover.)}

Bài học: một cười có thể tổn thương, nhưng sự khích lệ nhẹ nhàng sau đó có thể reparation được.
\textit{(The lesson: one laugh can wound, but gentle encouragement afterward can heal.)}
\end{truyen}

% Story 52
\begin{truyen}{52. Không Dấy Tay Trả Lời Câu Hỏi}{Never Raising Their Hand}
Em Bình thông minh, nhưng từ lớp 6, em không bao giờ dấy tay trong lớp.
\textit{(Binh was smart, but since 6th grade, he never raised his hand in class.)}
Tại sao?
\textit{(Why?)}

Vì một lần, em trả lời và bị thầy nói: ``Sai rồi, vì sao em không nghe kỹ?''
\textit{(Because once, he answered and the teacher said: ``Wrong, why didn't you listen carefully?'')}

Một từ đơn chỉ là phê bình, nhưng đối với em Bình nó là: ``Em ngu.''
\textit{(One phrase was just criticism, but for Binh it meant: ``You're stupid.'')}

Em Bình bắt đầu nói với bản thân: ``Nếu mình có thể sai, thì tốt hơn là đừng nói gì.''
\textit{(Binh started telling himself: ``If I can be wrong, it's better to say nothing.'')}

Sau này, kỳ thi cuối năm, những em khác trao đổi bài để học hỏi, nhưng Bình không.
\textit{(Later, at the end-of-year exam, other students discussed their answers to learn, but not Binh.)}
Em e rằng nếu nói, em sẽ bị phán xét.
\textit{(He was afraid if he spoke, he'd be judged.)}

Một ngày, thầy Hiền - thầy dạy khi Bình lên lớp 8 - đặc biệt trò chuyện với em:
\textit{(One day, Mr. Hien -- Binh's 8th grade teacher -- specifically talked to him:)}
``Mỗi câu trả lời sai là một lần học tập. Tại sao em lại sợ học tập?''
\textit{(``Every wrong answer is a learning opportunity. Why would you fear learning?'')}

Em Bình lần đầu tiên thấy, sự sợ hãi không cho, sai lầm là một phần của học hỏi.
\textit{(For the first time, Binh understood: fear doesn't protect, mistakes are part of learning.)}

Bài học: giáo viên có quyền lực để phá hoại hoặc xây dựng lại tự tin của học sinh.
\textit{(The lesson: teachers have power to destroy or rebuild student confidence.)}
\end{truyen}

% Story 53
\begin{truyen}{53. Chiếc Áo Đồng Phục Cũ}{The Worn-Out Uniform}
Em Loan lớp 6 có một chiếc áo đồng phục duy nhất.
\textit{(Loan in 6th grade had only one school uniform.)}
Em mặc nó ba ngày một lần, rửa nó ba đêm một lần để khi sáng, nó đã khô.
\textit{(She wore it three days, washed it three nights so by morning it was dry.)}

Những em khác có năm cái áo, mỗi cái bề mặt xanh lam sạch sẽ.
\textit{(Other students had five uniforms, each a crisp clean blue.)}
Áo của Loan bề mặt và bị phai màu từ nhiều lần giặt.
\textit{(Loan's was faded from many washings.)}

Một hôm thầy quân sự nhận xét: ``Áo em bê bối quá. Sao không mặc cái áo khác ?''
\textit{(One day, the military instructor remarked: ``Your uniform is shabby. Why don't you wear another one?'')}

Lớp cười.
\textit{(The class laughed.)}
Em Loan kéo áo lại, ngồi dự tiết học kế tiếp mà cảm thấy như một ai đó khác, không phải em Loan.
\textit{(Loan pulled her shirt tight, sat through the next class feeling like someone else, not herself.)}

Ngày hôm sau, em không đi học.
\textit{(The next day, she didn't go to school.)}

Bố mẹ hỏi tại sao.
\textit{(Her parents asked why.)}
Em không nói, chỉ nằm trên giường.
\textit{(She didn't answer, just lay in bed.)}

Một cô giáo chủ nhiệm tới thăm nhà.
\textit{(The homeroom teacher came to visit.)}
Khi biết sự thật, cô quay lại lớp và nói:
\textit{(When she learned the truth, she went back to class and said:)}

``Hiện tại, không phải ai cũng có từng chi có tất cả những gì họ muốn. Nhưng một lựa chọn khác là tri ơn những người mặc cùng một áo và cảm thấy như anh chị em.''
\textit{(``Right now, not everyone can have everything they want. But another choice is to honor those wearing the same uniform and feel like brothers and sisters.'')}

Em Loan quay lại với một cảm tưởng khác.
\textit{(Loan came back with a different feeling.)}

Bài học: xấu hổ không phải là điều kiện của em, nó là một lựa chọn cách nhìn.
\textit{(The lesson: shame isn't a condition, it's a choice of perspective.)}
\end{truyen}

% Story 54
\begin{truyen}{54. Khen Muộn Năm Sau}{Praise That Comes Too Late}
Em Cường lớp 6 học rất giỏi.
\textit{(Cuong in 6th grade was a top student.)}
Em là em đầu tiên dẫn lớp, luôn ghi chép nhanh, luôn xin bị kiểm tra.
\textit{(He was first in class, always took notes quickly, always asked for quizzes.)}

Nhưng trong quá trình từ lớp 6 tới lớp 7, không ai khen em.
\textit{(But from 6th to 7th grade, no one praised him.)}
Em làm tốt, nhưng không có ai nhận ra.
\textit{(He did well, but no one noticed.)}

Em bắt đầu nghĩ: ``Cái gì?''
\textit{(He started thinking: ``What for?'')}
``Nếu tôi làm tốt, nó cũng không có ý nghĩa gì.''
\textit{(``If I do well, it doesn't matter.'')}

Lớp 7, em bắt đầu bỏ học, sao chép, không tìm tòi.
\textit{(In 7th grade, he started skipping, copying, not looking things up.)}
Em từ học sinh giỏi thành học sinh yếu.
\textit{(He went from top student to weak student.)}

Lớp 8, thầy mới nhận thấy.
\textit{(In 8th grade, a new teacher noticed.)}
Thầy khen em rất kỳ lạ.
\textit{(The teacher praised him strangely.)}
``Em Cường có trí tuệ. Sao em lại bỏ lãng nó?''
\textit{(``You have intelligence, Cuong. Why have you abandoned it?'')}

Nhưng quá muộn rồi.
\textit{(But it was too late.)}
Em Cường cảm thấy như những khen ngợi bây giờ nó là tuyệt vọng, không phải động lực.
\textit{(Cuong felt like the praise now was despair, not motivation.)}

Bài học: khen ngợi sớm, khen ngợi cụ thể, khen ngợi có ý nghĩa.
\textit{(The lesson: praise early, praise specifically, praise meaningfully.)}
Khen muộn còn tồi tệ hơn không khen, vì nó để lại cảm giác bị lãng quên.
\textit{(Late praise is worse than no praise, because it leaves the feeling of being forgotten.)}
\end{truyen}

% Story 55
\begin{truyen}{55. Hy Vọng Vô Hy Vọng}{The Hopeless Label}
Lớp nào cũng có em được gọi bằng "bê bối" hoặc "vô vọng".
\textit{(Every class has a student called ``hopeless'' or ``useless''.)}
Em Tú lớp 8 là em như vậy.
\textit{(Tú in 8th grade was one.)}

``Thằng Tú thì từ kỳ 1 đến kỳ 2 không có gì thay đổi. Thằng này hình như bị vô năng vậy.''
\textit{(``Tú never changes from semester 1 to 2. This kid seems unable to learn.'')}
Đó là những gì những thầy nói với nhau.
\textit{(That's what the teachers said about him.)}

Em Tú nghe qua những bạn bè.
\textit{(Tú heard through friends.)}
Em hiểu rằng người lớn đã từ bỏ em, cơ hội cũng ít dần, và trên người em như bị dán chặt một cái nhãn: ``vô vọng.''
\textit{(He understood: adults had given up, he was out of chances, he was branded with a label: ``hopeless''.)}

Năm thứ ba, em Tú không đi học nữa.
\textit{(In his third year, Tú stopped going to school.)}

Nhưng cô giáo tiếng Anh -- cô Huế -- không tận tình từ bỏ.
\textit{(But his English teacher -- Ms. Hue -- didn't give up.)}
Cô đi thăm nhà em, nói: ``Tú có khả năng, chỉ là sợ.''
\textit{(She visited his home and said: ``You have potential, you're just afraid.'')}

``Lần này, hãy thử một lần nữa. Cô sẽ đi cùng với em.''
\textit{(``This time, try once more. I'll go with you.'')}

Em Tú trở lại.
\textit{(Tú came back.)}
Em không trở thành học sinh xuất sắc ngay lập tức, nhưng em không còn bị đè bẹp bởi cảm giác vô vọng nữa.
\textit{(Not suddenly brilliant, but no longer paralyzed by hopelessness.)}

Bài học: một nhãn dán có thể cứu người hoặc có thể giết.
\textit{(The lesson: a label can save someone or kill them.)}
\end{truyen}

% Story 56
\begin{truyen}{56. Nói Qua Bạn}{Speaking Through a Friend}
Em Hạnh có ý kiến tốt, nhưng sợ nói trực tiếp.
\textit{(Hanh had good ideas, but was afraid to say them directly.)}
Em nói với thầy nói: ``Thầy, em của chúng em muốn hỏi ...''
\textit{(She'd tell the teacher: ``Teacher, my friend wants to ask...'')}
Em không bao giờ nói: ``Tôi muốn.''
\textit{(She never said: ``I want.'')}

Một ngày, cô giáo nhận ra.
\textit{(One day, the teacher noticed.)}
Cô gọi em nói chuyện: ``Em nói giúp bạn, nhưng bạn em có biết cách nói không?''
\textit{(She called her over: ``You speak for your friend, but does she know how to speak for herself?'')}

Em Hạnh im lặng.
\textit{(Hanh was silent.)}

Cô tiếp: ``Em không phải yếu, em chỉ là sợ. Nhưng sợ là cái gì em có thể thay đổi được.''
\textit{(The teacher continued: ``You're not weak, you're just afraid. But fear is something you can change.'')}

Từ hôm đó, cô bắt em Hạnh nói trực tiếp.
\textit{(From that day on, the teacher made her speak directly.)}
Căng thẳng lần đầu.
\textit{(It was stressful at first.)}
Nhưng rồi em Hạnh nhận ra: khi cô giáo chỉ nghe em nói trực tiếp, em dần mạnh dạn hơn thật.
\textit{(But then Hanh discovered: the teacher would only listen to her speak directly, making her stronger.)}

Bài học: nếu em cho phép bị im lặng, thế giới sẽ từ bỏ lỗi em.
\textit{(The lesson: if you allow yourself to be silenced, the world forgets you exist.)}
\end{truyen}

% Story 57
\begin{truyen}{57. Sợ Một Bộ Môn Học}{The Subject You're Terrified Of}
Em Kiệt sợ Toán từ lớp 5.
\textit{(Kiet had been afraid of math since 5th grade.)}
Bất cứ khi nào thầy dạy Toán, trái tim em đập nhanh, tay em run.
\textit{(Whenever the math teacher taught, his heart raced, his hands shook.)}

Tại sao?
\textit{(Why?)}
Vì một lần, thầy Cương gọi em lên bảng.
\textit{(Because once, Teacher Cuong called him to the board.)}
Em không làm được bài, đứng bối rối trước lớp.
\textit{(He couldn't solve it, stood confused in front of the class.)}

Thầy cười và nói: ``Em thằng này không dành cho Toán, quên nó đi.''
\textit{(The teacher laughed and said: ``You're just not made for math, give it up.'')}

Em Kiệt quên không được.
\textit{(Kiet couldn't forget.)}
Những lần tiếp theo, em e dè, sợ hãi, sợ được gọi lên.
\textit{(Next times, he was anxious, terrified of being called up.)}

Lớp 8, thầy dạy Toán khác.
\textit{(In 8th grade, a different math teacher.)}
Thầy nhận thấy em Kiệt sợ hãi trong giờ học.
\textit{(The teacher noticed Kiet's fear.)}
Thầy không bắt em trả lời trước lớp.
\textit{(He didn't make him answer in front of the class.)}
Thầy hỏi em riêng, từ từ, với lòng kiên nhẫn.
\textit{(He asked him privately, slowly, with patience.)}

Rồi em Kiệt bắt đầu hiểu, những bài Toán không phải là họa nạn.
\textit{(Then Kiet began to understand, math wasn't a disaster.)}
Nó là ngôn ngữ.
\textit{(It was language.)}

Bài học: sợ không phải là bản chất của bạn, nó là một ký ức cũ chưa chết.
\textit{(The lesson: fear isn't who you are, it's an old memory not yet laid to rest.)}
\end{truyen}

% Story 58
\begin{truyen}{58. Bị Chỉ Trích Ngoài Hành Lang}{Corrected in the Hallway}
Em Huy làm sai một bước trong kế hoạch CLB.
\textit{(Huy made one mistake in the club plan.)}
Thầy Hùng gặp em trên hành lang và nói to: ``Em này, làm gì mà sai hoàn toàn vậy!''
\textit{(Teacher Hung met him in the hallway and said loudly: ``You, how could you get it completely wrong!'')}

Cả hành lang nghe thấy.
\textit{(The whole hallway heard.)}
Em Huy cảm thấy tất cả mọi người đều nhìn em, đánh giá em, kết luận: ``thằng Huy là kẻ vô dụng.''
\textit{(Huy felt everyone looking at him, judging him, concluding: ``Huy is useless.'')}

Một tuần sau, em Huy sợ đi học.
\textit{(A week later, Huy was afraid to go to school.)}
Em không muốn gặp thầy Hùng, không muốn gặp bạn bè.
\textit{(He didn't want to see Teacher Hung, didn't want to see friends.)}

Sau đó, một lần, thầy Hùng hiểu ra.
\textit{(Later, Teacher Hung realized.)}
Thầy gọi em ra riêng: ``Em, thầy không có ý chỉ trích em nơi công cộng. Thầy xin lỗi. Lần tới, thầy sẽ nói chuyện với em riêng.''
\textit{(He called Huy privately: ``I shouldn't have criticized you publicly. I'm sorry. Next time, we'll talk privately.'')}

Em Huy quay lại.
\textit{(Huy came back.)}

Bài học: cách chỉ lỗi quan trọng hơn nội dung chỉ lỗi.
\textit{(The lesson: how you point out mistakes matters more than the mistake itself.)}
Chỉ lỗi nơi công cộng có thể để lại một vết sẹo suốt đời.
\textit{(Public criticism can leave a scar for life.)}
\end{truyen}

% Story 59
\begin{truyen}{59. Cuộc Gặp Sau Khi Sai Lầm}{The After-Mistake Meeting}
Em Sang làm bài tập nhóm, nhưng không làm phần của em, bỏ hết cho bạn.
\textit{(Sang didn't do his group project part, left everything to his friends.)}
Bạn em phát hiện, nói với cô.
\textit{(His friends discovered it and told the teacher.)}

Cô gọi em vào văn phòng.
\textit{(The teacher called him to the office.)}

Em chắc chắn là sẽ bị la mắng, bị kiểm điểm, bị giáng chỉ tiêu.
\textit{(He was sure he'd be scolded, criticized, and penalized.)}

Nhưng cô hỏi: ``Em tại sao không làm?''
\textit{(But she asked: ``Why didn't you do your part?'')}
Em Sang rụt rè nói: ``em không hiểu bài.''
\textit{(Sang hesitantly said: ``I didn't understand the lesson.'')}

Cô nói: ``Vậy em đáng lẽ phải nói với cô hoặc bạn em, chứ không phải bỏ hết.''
\textit{(She said: ``You should have told me or your friends, not just given up.'')}
``Lần sau, em hãy nói trước khi bỏ cuộc.''
\textit{(``Next time, speak up before you quit.'')}

Em Sang cảm thấy bị lắng nghe chứ không bị lên án.
\textit{(Sang felt heard, not condemned.)}

Bài học: sau một sai lầm, cái em cần không phải là hình phạt mà là sự hiểu biết và cơ hội sửa sai.
\textit{(The lesson: after a mistake, what you need isn't punishment but understanding and a chance to correct it.)}
\end{truyen}

% Story 60
\begin{truyen}{60. Ngã Về Phía Trước}{Falling Forward}
Em Dạo bị phạt từ lớp sau nhiều lần bỏ học.
\textit{(Dao was expelled from school after many absences.)}
Em cảm thấy như một kẻ thất bại.
\textit{(He felt like a failure.)}

Nhưng bố mẹ em không buộc tội.
\textit{(But his parents didn't blame him.)}
Họ hỏi: ``Em muốn làm gì bây giờ, con?''
\textit{(They asked: ``What do you want to do now, son?'')}

Em Dao bối rối.
\textit{(Dao was confused.)}
``Con không biết, ba mẹ.''
\textit{(``I don't know, Mom and Dad.'')}

``Thôi, thử học thủ công tạo hình. Nếu em không thích, chúng ta sẽ tìm cái khác.''
\textit{(``Alright, try vocational training. If you don't like it, we'll find something else.'')}

Sau một năm, em Dao trở thành thợ thủ công có tay nghề.
\textit{(After a year, Dao became a skilled craftsman.)}
Em không trở thành giáo sư, nhưng em trở thành người có tay nghề, tự tin.
\textit{(He didn't become a professor, but he became skilled, confident.)}

Một ngày, em về thăm cô giáo cũ.
\textit{(One day, he visited his old teacher.)}
Cô giáo hỏi: ``Em vừa bị phạt từ lớp, em có hối hận không?''
\textit{(The teacher asked: ``You were expelled, do you regret it?'')}

Em Dao nói: ``Cô, em không hối hận cái lựa chọn của mình. Em chỉ hối hận là em đã dùng thời gian quá lâu để phát hiện: em không phải là suy sụp, em chỉ là đi một con đường khác.''
\textit{(Dao said: ``Teacher, I don't regret my choice. I just regret it took me so long to discover: I'm not a failure, I just went a different path.'')}

Bài học: xấu hổ là những khoảnh khắc khi em hiểu mình con người ra sao.
\textit{(The lesson: shame is when you discover who you really are.)}
Nó chỉ trở thành thất bại khi em chọn không đứng dậy.
\textit{(It only becomes failure if you choose not to rise.)}
\end{truyen}

\begin{tuhocnhanh}{Câu Hỏi Tự Học Chương 6}{Self-Study Questions Chapter 6}

\textbf{Câu 1:} Em đã từng bị xấu hổ vì cái gì?
\textbf{Question 1:} Have you ever felt ashamed of something?

\textbf{Câu 2:} Xấu hổ ấy, em đó có làm em sợ hãi hoặc làm em thay đổi hành động?
\textbf{Question 2:} Did that moment of shame make you fearful or change how you act?

\textbf{Câu 3:} Nếu nhìn lại, em có thấy xấu hổ ấy là cơ hội chứ không phải là kết luận cuối cùng không?
\textbf{Question 3:} Looking back, do you see that shame as an opportunity, not a final verdict?

\textbf{Câu 4:} Em muốn những người khác nhắc em khi em sai lầm bằng cách nào?
\textbf{Question 4:} How do you want others to remind you when you make mistakes?

\end{tuhocnhanh}
